% Options for packages loaded elsewhere
\PassOptionsToPackage{unicode}{hyperref}
\PassOptionsToPackage{hyphens}{url}
\PassOptionsToPackage{dvipsnames,svgnames,x11names}{xcolor}
\documentclass[
  10pt,
  fontset=fandol]{ctexart}
\usepackage{xcolor}
\usepackage[margin=16mm]{geometry}
\usepackage{amsmath,amssymb}
\setcounter{secnumdepth}{-\maxdimen} % remove section numbering
\usepackage{iftex}
\ifPDFTeX
  \usepackage[T1]{fontenc}
  \usepackage[utf8]{inputenc}
  \usepackage{textcomp} % provide euro and other symbols
\else % if luatex or xetex
  \usepackage{unicode-math} % this also loads fontspec
  \defaultfontfeatures{Scale=MatchLowercase}
  \defaultfontfeatures[\rmfamily]{Ligatures=TeX,Scale=1}
\fi
\usepackage{lmodern}
\ifPDFTeX\else
  % xetex/luatex font selection
\fi
% Use upquote if available, for straight quotes in verbatim environments
\IfFileExists{upquote.sty}{\usepackage{upquote}}{}
\IfFileExists{microtype.sty}{% use microtype if available
  \usepackage[]{microtype}
  \UseMicrotypeSet[protrusion]{basicmath} % disable protrusion for tt fonts
}{}
\makeatletter
\@ifundefined{KOMAClassName}{% if non-KOMA class
  \IfFileExists{parskip.sty}{%
    \usepackage{parskip}
  }{% else
    \setlength{\parindent}{0pt}
    \setlength{\parskip}{6pt plus 2pt minus 1pt}}
}{% if KOMA class
  \KOMAoptions{parskip=half}}
\makeatother
\setlength{\emergencystretch}{3em} % prevent overfull lines
\providecommand{\tightlist}{%
  \setlength{\itemsep}{0pt}\setlength{\parskip}{0pt}}
\usepackage{amsmath,amssymb,mathtools,needspace}
\xeCJKDeclareCharClass{CJK}{"2032,"0400->"04FF,"2160->"216B,"2460->"2473,"25A0,"25C7,"25CB,"25CF}
\IfFontExistsTF{Arial Unicode MS}{\xeCJKsetup{AutoFallBack=true}\setCJKfallbackfamilyfont{\CJKrmdefault}{Arial Unicode MS}}{}
\setcounter{secnumdepth}{0}
\setlength{\emergencystretch}{2em}
\allowdisplaybreaks
\usepackage{bookmark}
\IfFileExists{xurl.sty}{\usepackage{xurl}}{} % add URL line breaks if available
\urlstyle{same}
\hypersetup{
  pdftitle={变步长的 Runge-Kutta 方法},
  colorlinks=true,
  linkcolor={Maroon},
  filecolor={Maroon},
  citecolor={Blue},
  urlcolor={Blue},
  pdfcreator={LaTeX via pandoc}}

\title{变步长的 Runge-Kutta 方法}
\author{}
\date{}

\begin{document}
\maketitle

\subsubsection{5.3.6 变步长的 Runge-Kutta 方法}\label{ux53d8ux6b65ux957fux7684-runge-kutta-ux65b9ux6cd5}

单从每一步看，步长越小，截断误差就越小；但随着步长的缩小，在一定求解范围内所要完成的步数就增加了。步数的增加不但引起计算量的增大，而且可能导致舍入误差的严重积累。因此同积分的数值计算一样，微分方程的数值解法也有个选择步长的问题。

在选择步长时，需要考虑两个问题：

（1）怎样衡量和检验计算结果的精度？

（2）如何依据所获得的精度处理步长？

考察经典的四阶 Runge-Kutta 格式 (5.3.13)。从节点 \(x_n\) 出发，先以 \(h\) 为步长求出一个近似值，作为 \(y_{n+1}^{(h)}\)，由于格式的局部截断误差为 \(O(h^5)\)，故有

\[
y(x_{n+1})-y_{n+1}^{(h)}\approx Ch^5;\tag{5.3.14}
\]

然后将步长折半，即取 \(\dfrac{h}{2}\) 为步长从 \(x_n\) 跨两步到 \(x_{n+1}\)，再求得一个近似值 \(y_{n+1}^{(h/2)}\)，每跨一步的截断误差是 \(C\left(\dfrac{h}{2}\right)^5\)，因此有

\[
y(x_{n+1})-y_{n+1}^{(h/2)}=2C\left(\frac{h}{2}\right)^5;\tag{5.3.15}
\]

比较式 (5.3.14) 和式 (5.3.15) 可以看到，步长折半后，误差大约减小到 \(\dfrac{1}{16}\)，即有

\[
\frac{y(x_{n+1})-y_{n+1}^{(h/2)}}{y(x_{n+1})-y_{n+1}^{(h)}}\approx\frac{1}{16}.
\]

由此易得下列事后估计式

\[
y(x_{n+1})-y_{n+1}^{(h/2)}\approx\frac{1}{15}(y_{n+1}^{(h/2)}-y_{n+1}^{(h)}).
\]

\href{../../source-images/pdf-133.jpeg}{原页}

这样，可以通过检查步长折半前后两次计算结果的偏差

\[
\Delta=|y_{n+1}^{(h/2)}-y_{n+1}^{(h)}|
\]

来判定所选的步长是否合适。具体地说，将分以下两种情况处理：

（1）对于给定的精度 \(\varepsilon\)，如果 \(\Delta>\varepsilon\)，应反复将步长折半进行计算，直至 \(\Delta<\varepsilon\) 为止，这时取最终得到的 \(y_{n+1}^{(h/2)}\) 作为结果；

（2）如果 \(\Delta<\varepsilon\)，则反复将步长加倍，直至 \(\Delta>\varepsilon\) 为止，这时再将步长折半一次，就得到所要的结果。

这种通过加倍或折半处理步长的方法称为\textbf{变步长方法}。表面上看，为了选择步长，每一步的计算量增加了，但从总体考虑往往是合算的。

\end{document}
